\section{Punkty idealne i ultraidealne} % lang-pl
\label{sekcja_punkty_idealne}

Na płaszczyźnie euklidesowej dwie proste albo się przecinają, albo są równoległe; na hiperbolicznej dochodzi trzeci przypadek -- proste hiperrównoległe. % lang-pl
W tej sekcji zobaczymy, jak Hilbert opisze ten trzeci przypadek i jak uzupełnić płaszczyznę fikcyjnymi punktami, tak jak w geometrii euklidesowej dokłada się punkty w nieskończoności. % lang-pl

\begin{theorem}
[Hilberta] % lang-pl
	Dwie proste hiperrównoległe mają dokładnie jedną wspólną prostopadłą. % lang-pl
\end{theorem}

Dowód Hilberta jest konstruktywny: Eves \cite[s. 305--306]{eves1_1972}. % lang-pl
Eves zwraca uwagę, że konstrukcje w geometrii Łobaczewskiego (równoległych przez dany punkt, wspólnej równoległej do dwóch przecinających się prostych) bynajmniej nie są oczywiste \cite[s. 306]{eves1_1972}. % lang-pl
\index[persons]{Hilbert, David}%

\begin{definition}
[punkt idealny, punkt ultraidealny] % lang-pl
\index{punkt!idealny} % lang-pl
\index{punkt!ultraidealny} % lang-pl
	Każdej rodzinie półprostych równoległych przypisujemy wspólny punkt idealny. % lang-pl
	Każdej rodzinie prostych hiperrównoległych, prostopadłych do jednej prostej, przypisujemy wspólny punkt ultraidealny. % lang-pl
\end{definition}

Eves \cite[s. 306]{eves1_1972}. % lang-pl

Wszystko to da się zobaczyć na jednym rysunku: zwykłe punkty płaszczyzny Łobaczewskiego to punkty wewnątrz ustalonego koła $K$, punkty idealne leżą na okręgu $K$, ultraidealne na zewnątrz (łącznie z nieskończonymi), a proste to proste płaszczyzny rozszerzonej przecinające wnętrze koła; prosta stowarzyszona z punktem ultraidealnym to jego biegunowa względem $K$ \cite[s. 306--307]{eves1_1972}. % lang-pl
Arthur Cayley nazwie okrąg $K$ \emph{absolutem} płaszczyzny (zobacz sekcję \ref{sekcja_beltrami_cayley_klein}). % lang-pl
Eves widzi w tym na razie tylko wygodny rysunek, który pozwala zapamiętać wiele faktów; dopiero po przygotowaniu tła posłuży on do pokazania związku z geometrią rzutową \cite[s. 307]{eves1_1972}. % lang-pl
\index{absolut} % lang-pl

W tym obrazie każde dwie proste przecinają się w punkcie zwykłym, idealnym lub ultraidealnym, ale nie każde dwa punkty wyznaczają prostą \cite[s. 307]{eves1_1972}. % lang-pl

\begin{theorem}
	Symetralne boków dowolnego trójkąta przecinają się w jednym punkcie -- zwykłym, idealnym albo ultraidealnym. % lang-pl
\end{theorem}

Dowód (rozbity na trzy przypadki): Eves \cite[s. 308--311]{eves1_1972}. % lang-pl

Zadania: sprawdzenie trzynastu faktów w modelu absolutu -- Eves \cite[s. 307--308]{eves1_1972}. % lang-pl
% Źródła: Hartshorne s. 312--317 (§34, równoległość graniczna), 388 nn. (§41, końce).

